% §15.3 Space and Time Analysis
%
% Monograph register: rank table (n*r), perfect-hash storage (inherited
% from ch14), total, prop:bphf-space, expected O(n) build, O(1) query,
% comparison table at n=10^6 with the eps=1/n crossover. Trims the
% repeated "asymmetry is not a defect" and the doubled perfect-hash
% inheritance. Cross-refs def:bloom-filter, thm:entropy-floor, def:bhf,
% thm:bhf-optimality, sec:space-analysis, sec:construction-time, def:bphf,
% sec:bphf-construction, alg:bphf-build-query, eq:bphf-r, eq:bphf-build,
% eq:bphf-fnr, sec:ranked-search. DEFINED: sec:space-time-analysis,
% prop:bphf-space, eq:bphf-rank-bits, eq:bphf-total-bits,
% eq:bphf-per-element, tab:bhf-bphf-bits.

\section{Space and Time Analysis}
\label{sec:space-time-analysis}

The BPHF stores the perfect hash $\phi$ and the rank table $R$ of width
$r = \lceil \log_2 n \rceil$. This section totals the bits, gives the
build and query times the construction inherits from its chapter~14
ancestor, and tabulates the trade against the BHF. The headline is
\Cref{prop:bphf-space}: the per-element cost is $\lceil \log_2 n \rceil +
o(1)$, with the perfect-hash overhead vanishing as in
\Cref{sec:space-analysis}. The build is expected $O(n)$ and queries
$O(1)$, matching the BHF on the time axes. The structural difference is
on the space axis: the BPHF's bits-per-element depends on the set size
$n$, the BHF's on the target rate $\fprate$, so the two families are not
substitutes parameter-for-parameter.

\subsection*{Rank table storage}

The table $R : [n] \to \{0, 1\}^r$ holds $n$ slots of $r$ bits, so
\begin{equation}
\label{eq:bphf-rank-bits}
  L_R \;=\; n \cdot r \;=\; n \cdot \lceil \log_2 n \rceil
\end{equation}
bits, growing as $n \log_2 n$, the same information-theoretic cost a
sorted lookup index pays to distinguish $n$ ranks. At $n = 10^6$ the
table is $20 \cdot 10^6$ bits, about $2.5$ megabytes; at $n = 10^9$ it is
$30 \cdot 10^9$ bits, about $3.75$ gigabytes. The integer rounding in
$r = \lceil \log_2 n \rceil$ adds between $0$ and $1$ bits per element
over $\log_2 n$, zero at powers of two and averaging $0.5$, the rounding
profile \Cref{sec:space-analysis} showed for fingerprints; fractional-bit
rank encodings are not worth the query-time complexity, the same
conclusion chapter~14 reached. The table is a flat array addressed by
$\phi$, with no padding, the densest representation of one rank per
element.

\subsection*{Perfect-hash storage}

The perfect hash is a property of the set, not of the values stored at
each slot, so the BPHF inherits the BHF's overhead unchanged: the
Belazzougui-Botelho-Dietzfelbinger family gives
\[
  L_\phi \;=\; O\bigl( n \log \log \log n / \log n \bigr)
\]
bits, $o(n)$ and therefore vanishing relative to the rank table's
$n \log_2 n$, at practical constants of $2$ to $3$ bits per element at
$n = 10^6$ and below that at larger $n$. The chapter~14 build pass
populated $F[\phi(a)] = g(a)$; the chapter~15 pass populates
$R[\phi(a)] = \rho(a)$. Only the per-slot value differs; the
perfect-hash representation, slot count, and addressing are identical, so
the chapter does not redo the derivation.

\subsection*{Total bits per element}

Combining,
\begin{equation}
\label{eq:bphf-total-bits}
  L_{\text{BPHF}}
  \;=\; L_R + L_\phi
  \;=\; n \cdot \lceil \log_2 n \rceil
  \;+\; O\bigl( n \log \log \log n / \log n \bigr),
\end{equation}
and dividing by $n$,
\begin{equation}
\label{eq:bphf-per-element}
  L_{\text{BPHF}} / n
  \;=\; \lceil \log_2 n \rceil
  \;+\; O(\log \log \log n / \log n).
\end{equation}

\begin{proposition}[BPHF space]
\label{prop:bphf-space}
The Bernoulli Perfect Hash Filter of \Cref{def:bphf} on a designated set
$A$ with $|A| = n$ uses
\[
  L_{\text{BPHF}} / n
  \;=\; \lceil \log_2 n \rceil
  \;+\; o(1)
  \;=\; \log_2 n \;+\; O(1)
\]
bits per element, the rank table dominating at $\lceil \log_2 n \rceil$
bits and the perfect-hash overhead a vanishing sublinear correction.
\end{proposition}

The proposition has the form of \Cref{thm:bhf-optimality} with one
difference: the leading term is controlled by the set size $n$, not the
target rate $\fprate$. The two budgets
\[
  L_{\text{BHF}} / n = \lceil -\log_2 \fprate \rceil + o(1),
  \qquad
  L_{\text{BPHF}} / n = \lceil \log_2 n \rceil + o(1)
\]
are functions of different inputs. The BHF spends bits at a rate the
designer sets through $\fprate$; the BPHF at a rate the set size
dictates. Neither is a defect: each budget is the minimum its guarantee
requires. The BHF pays $\lceil -\log_2 \fprate \rceil$ bits to make a
false-positive coincidence rare; the BPHF pays $\lceil \log_2 n \rceil$
bits so the rank assignment can name all $n$ members apart and certify a
witness. At full resolution the BPHF is exact (the $\fnrate = 0$ corner);
the $k$-perfect dial of \Cref{sec:bphf-construction} shrinks the slot
count to $m < n$ and trades $\fnrate \approx 1 - m/n$ for the saved
space.

\subsection*{Build and query time}

The build writes one rank per element in a one-pass $O(n)$ stage plus the
expected-$O(n)$ perfect-hash construction of \Cref{sec:construction-time}
(the same Fredman-Komlos-Szemeredi, Hagerup-Tholey, MWHC, BBHash, and BDZ
families), dominated by the latter; build throughputs and times match the
BHF's. The rank-fill is embarrassingly parallel (distinct slots by
bijectivity), running after the perfect-hash stage in the same two-stage
pipeline. The only build-time difference from the BHF is computing a rank
$\rho(a)$ rather than a fingerprint $g(a)$ per element, both $O(1)$.

The query of \Cref{alg:bphf-build-query} runs in $O(1)$: one
perfect-hash evaluation, one rank evaluation, one memory read, and an
$r$-bit comparison (at machine-word speed for any $n$ up to $2^{64}$).
The cycle count is essentially the BHF's, differing only in which
function is the second evaluation and what the table read returns.

\subsection*{The bit-cost picture at \texorpdfstring{$n = 10^6$}{n = 10\^{}6}}

Fix $n = 10^6$. The BPHF's per-element cost is
$\lceil \log_2 10^6 \rceil = 20$ bits, independent of any error target;
the BHF's is $\lceil -\log_2 \fprate \rceil$. \Cref{tab:bhf-bphf-bits}
puts them side by side.

\begin{table}[h]
\centering
\caption{Bits per element at $n = 10^6$. BHF column is $\lceil -\log_2
\fprate \rceil$ (\Cref{thm:bhf-optimality}); BPHF column is $\lceil
\log_2 n \rceil = 20$ (\Cref{prop:bphf-space}), constant since the BPHF
guarantees $\fprate = 0$. Both omit the sublinear perfect-hash overhead.}
\label{tab:bhf-bphf-bits}
\begin{tabular}{l|cc}
\hline
  $\fprate$  & BHF & BPHF \\
\hline
  $10^{-2}$  & $7$  & $20$ \\
  $10^{-3}$  & $10$ & $20$ \\
  $10^{-6}$  & $20$ & $20$ \\
\hline
\end{tabular}
\end{table}

The two filters cross over at $\fprate = 1/n$, where $\lceil -\log_2
\fprate \rceil = \lceil \log_2 n \rceil$: at $n = 10^6$ this is $\fprate =
2^{-20} \approx 10^{-6}$. Above the crossover (looser targets, $\fprate >
1/n$) the BHF is cheaper; below it (stricter targets) the BPHF is
actually cheaper at the same $n$. The two regimes are the two operating
modes, the BHF where the application tolerates some false positives for
cheap storage, the BPHF where it cannot tolerate any. The BPHF's extra
bits in the loose-target regime are not wasted: they buy a guarantee the
BHF cannot deliver at any bit budget. A BHF at $\fprate = 10^{-2}$ admits
roughly one false positive per hundred non-member queries no matter how
many bits it spends; a BPHF at the same $n$ admits zero. The $13$ extra
bits per element are the price of converting a $10^{-2}$ probability into
a certainty on the consequential axis, the conversion
\Cref{sec:ranked-search} shows paying for itself in top-$k$ retrieval.
